\section{Environment Resolution}
\label{sec:environment}

A key advancement that \karonte suggested is modeling multi-binary interactions through analyzing Inter Process Communications (IPCs).
During multi-binary data-flow analysis, \karonte propagates constraints from source binaries to sink binaries, including environment information extracted along the way.
The environment resolution in \mango employs a similar tactic.
However, \mango does not perform symbolic exploration or propagate constraints across binaries that interact through IPCs.
Instead, \mango directly links setter values and getter variables during data-flow analysis, which achieves \emph{cross-binary} vulnerability sink resolution.
Our solution is more scalable than \karonte because it avoids path explosion issues in symbolic exploration and runtime overhead involved in constraint solving.
\mango also applies the same environment resolution approach to front-end keywords analysis that \satc pioneered.

% Instead of propogating constraints we run our normal data-flow analysis, but with many more restrictions.
% We do not allow analysis of sub-functions and limit the allowable call-depth to 2.
% The goal of this environment analysis is to not only capture all keys for nvram and environment variables, but also all values in setter functions such as setenv or set\_nvram.
% We then create a firmware-level dictionary of these values.
% During our normal data-flow analysis, if one of the IPC getter functions is called, we make an effort to look-up the data in the dictionary.
% If we marked the value in the dictionary as fully resolvable then we pass the resolved values to our getter function handler.
% If the function has an entry in the dictionary but the data is marked as unresolved, then we similarly make the return of the getter function an unresolved value.
% Finally, if there is no entry for the key then we resolve the value to a pre-determined constant.
% Our intuition for this resolution is that if the key is not used to set a value in any other binary in the system then it is very likely a value pre-determined by the firmware itself.
